\documentclass[12pt,a4paper]{article}
\usepackage[utf8]{inputenc}
\usepackage[T1]{fontenc}
\usepackage{amsmath}
\usepackage{amssymb}
\usepackage{graphicx}
\usepackage[spanish]{babel}
\usepackage[a4paper, margin=2.5cm]{geometry}
\usepackage[european]{circuitikz}
\usetikzlibrary{babel}

\title{Análisis de Tensión por Superposición}
\author{Gemini Code Assist}
\date{\today}

\begin{document}

\maketitle

\section{Descripción del Circuito}

El circuito consta de dos ramas conectadas en paralelo entre un nodo 'A' y tierra (GND). Cada rama contiene una fuente de tensión en serie con una resistencia. El objetivo es determinar la fórmula para la tensión $V_A$ utilizando el teorema de superposición.

\begin{center}
    \begin{circuitikz}[scale=1.2, transform shape]
        \draw (0,0) node[ground] (gnd) {};
        \draw (0,3) node[circ, label=above:A] (A) {};
        
        % Rama 1: Fuente E1 y R1
        \draw (A) to[R, l=$R_1$] (-3,3) to[V, v<=$E_1$] (-3,0) -- (gnd);
        
        % Rama 2: Fuente E2 y R2
        \draw (A) to[R, l=$R_2$] (3,3) to[V, v<=$E_2$] (3,0) -- (gnd);
    \end{circuitikz}
\end{center}

\section{Aplicación del Teorema de Superposición}

El teorema de superposición nos permite calcular la contribución de cada fuente independiente por separado y luego sumar los resultados.

\subsection{Paso 1: Contribución de la fuente $E_1$}

Apagamos la fuente $E_2$ (la reemplazamos por un cortocircuito). El circuito se convierte en un divisor de tensión. La tensión en el nodo A, que llamaremos $V_{A1}$, es la tensión que cae en la resistencia $R_2$.

\begin{center}
    \begin{circuitikz}[scale=1.2, transform shape]
        \draw (0,0) node[ground] (gnd) {};
        \draw (0,3) node[circ, label=above:A] (A) {};
        \draw (A) to[R, l=$R_1$] (-3,3) to[V, v<=$E_1$] (-3,0) -- (gnd);
        \draw (A) to[R, l=$R_2$] (3,3) to[short] (3,0) -- (gnd);
    \end{circuitikz}
\end{center}

La tensión $V_{A1}$ se calcula con la fórmula del divisor de tensión:
\begin{equation}
    V_{A1} = E_1 \cdot \frac{R_2}{R_1 + R_2}
\end{equation}

\subsection{Paso 2: Contribución de la fuente $E_2$}

Ahora, apagamos la fuente $E_1$ (la reemplazamos por un cortocircuito) y calculamos la contribución de $E_2$. De nuevo, tenemos un divisor de tensión. La tensión en el nodo A, $V_{A2}$, es la que cae en la resistencia $R_1$.

\begin{center}
    \begin{circuitikz}[scale=1.2, transform shape]
        \draw (0,0) node[ground] (gnd) {};
        \draw (0,3) node[circ, label=above:A] (A) {};
        \draw (A) to[R, l=$R_1$] (-3,3) to[short] (-3,0) -- (gnd);
        \draw (A) to[R, l=$R_2$] (3,3) to[V, v<=$E_2$] (3,0) -- (gnd);
    \end{circuitikz}
\end{center}

La tensión $V_{A2}$ se calcula de forma análoga:
\begin{equation}
    V_{A2} = E_2 \cdot \frac{R_1}{R_1 + R_2}
\end{equation}

\subsection{Paso 3: Tensión Total en el Nodo A}

La tensión total $V_A$ es la suma de las contribuciones individuales:
\begin{equation}
    V_A = V_{A1} + V_{A2} = E_1 \frac{R_2}{R_1 + R_2} + E_2 \frac{R_1}{R_1 + R_2}
\end{equation}
Simplificando la expresión, obtenemos la fórmula final para la tensión en el nodo A:
\begin{equation}
    V_A = \frac{E_1 R_2 + E_2 R_1}{R_1 + R_2}
\end{equation}

\textit{Nota: Esta es la fórmula clásica del Teorema de Millman para dos ramas con fuentes de tensión.}

\end{document}